\subsection*{Sprint 5: Monitorización y métricas}
\addcontentsline{toc}{subsection}{Sprint 5: Monitorización y métricas}

El quinto sprint aborda la \textbf{monitorización} del sistema, esto es, la
capacidad de observar desde fuera qué está haciendo \textbf{Muncher} mientras
se ejecuta. Hasta el sprint anterior la aplicación se limitaba a una interfaz
servida como contenido estático, cuyo comportamiento se agota en el navegador
de quien la visita y puede comprobarse abriéndola. La incorporación de la API
ha cambiado esa situación: el sistema realiza ahora procesamiento en el lado
del servidor, ejecuta código en un entorno al que nadie asiste y lo hace a
través de varios componentes de nube que deben entenderse entre sí. Lo que
ocurre en esas ejecuciones no es observable salvo que el propio sistema lo
registre, de modo que la monitorización deja de ser una mejora deseable para
convertirse en la única vía de saber cómo se comporta el sistema.

Igual que ocurrió con los controles de calidad del primer sprint y con la
internacionalización del tercero, esta tarea se aborda en una fase temprana de
manera deliberada. Introducir la monitorización cuando el sistema es todavía
pequeño permite que cada funcionalidad nueva nazca ya instrumentada, en lugar
de tener que recorrer más adelante un código extenso para añadir la
instrumentación que se omitió. Además, los problemas que la monitorización
saca a la luz —tiempos de respuesta que se degradan, llamadas entre servicios
que fallan de forma intermitente, permisos mal ajustados entre componentes de
la nube— resultan mucho más baratos de corregir cuanto menos código y menos
infraestructura dependen de la parte afectada.

Anticipar esta tarea tiene un segundo efecto sobre la naturaleza de los errores
que el proyecto puede detectar. Un fallo de rendimiento o una interconexión
defectuosa entre servicios de la nube no se manifiestan necesariamente como una
excepción: el sistema sigue respondiendo, pero lo hace con lentitud, con
reintentos o de forma inconsistente. Sin datos de ejecución, esta clase de
comportamiento solo se descubre cuando alguien lo padece y sabe describirlo, lo
que en las etapas siguientes del proyecto —con la persistencia de datos y el
procesamiento de recetas todavía por incorporar— supondría depurar a ciegas
sobre una superficie cada vez mayor.

El propósito último del sprint es, por tanto, disponer de una visión clara del
comportamiento del sistema: saber qué peticiones recibe la API, cuánto tarda en
atenderlas, con qué frecuencia falla y en qué punto se produce el fallo. Esa
visión sirve tanto para diagnosticar incidencias concretas como para
fundamentar en datos las decisiones técnicas de los sprints posteriores, que
hasta ahora se han apoyado en el análisis previo y no en medidas tomadas sobre
el sistema en funcionamiento.

\subsubsection*{Objetivos iniciales}

\begin{itemize}
	\item Diseñar el sistema de monitorización, definiendo tanto los datos a recoger
	      como el stack de tecnologías que los procesará y almacenará.
	\item Implementar la instrumentación de la API para registrar las métricas definidas.
	\item Configurar la infraestructura de monitorización en la nube, incluyendo la
	      creación de los recursos necesarios para almacenar y visualizar las métricas.
	\item Desplegar la infraestructura de monitorización y la API instrumentada en el entorno
	      de desarrollo.
	\item Verificar que las métricas se recogen correctamente y que los paneles de
	      visualización muestran la información esperada.
	\item Desplegar la infraestructura de monitorización y la API instrumentada en el entorno de
	      producción.
\end{itemize}

\subsubsection*{Historias de usuario completadas}

\begin{center}
	\begin{tabularx}{\textwidth}{|l|X|c|}
		\hline
		\textbf{Identificador} & \textbf{Nombre}                                 & \textbf{Puntos de historia} \\
		\hline
		\usref{US-DE-04}       & Monitorización de la infraestructura en la nube & 5                           \\
		\hline
		\usref{US-DE-05}       & Instrumentación de la API                       & 3                           \\
		\hline
	\end{tabularx}
\end{center}

\subsubsection*{Retrospectiva}

El quinto sprint concluyó con el sistema observable en sus dos entornos: los
\textit{logs} de las funciones Lambda se recogen en CloudWatch, la API genera
trazas en X-Ray que desglosan cada petición por operación y el entorno de
producción dispone de un panel que reúne las métricas de la función y de la
distribución. Todo ello se declaró en la misma plantilla de infraestructura que
el resto del sistema, de modo que la monitorización se crea, se actualiza y se
destruye junto con el entorno al que pertenece.

La dificultad principal del sprint residió en el desglose de las trazas por
operación que exige \reqref{RNF-MO-03}. El rastreo activo de Lambda, configurado
en primer lugar, registra cada invocación y separa el arranque en frío de la
ejecución, pero trata esta última como un bloque indivisible. La vía habitual
para subdividirla, el SDK de X-Ray, había entrado en modo de mantenimiento y
tiene fecha de fin de soporte, por lo que adoptarlo habría supuesto programar
una migración a pocos meses vista. Se optó en su lugar por
\textbf{OpenTelemetry}, a través del \textit{layer} que AWS publica para Lambda.
La actualización de \textbf{FastAPI} a la versión 0.142, que emite de forma
nativa un \textit{span} por petición y otro por cada fase de su procesamiento,
permitió que la API quedara instrumentada sin una sola línea de código propio, y
que en local, donde no existe el SDK, la aplicación se comporte exactamente
igual que antes.

En cuanto a los requerimientos, el sprint satisfizo \reqref{RNF-MO-01} y
\reqref{RNF-MO-03}, y de forma parcial \reqref{RNF-MO-04}, ya que el recuento de
invocaciones mide peticiones atendidas y no usuarios activos.
\reqref{RNF-MO-02} quedó sin satisfacer de manera deliberada: ningún mecanismo
limita por sí solo el número de avisos que un sistema de alertas emite, y
contenerlo dentro de la capa gratuita habría exigido añadir infraestructura
capaz de fallar a su vez. Quedan asimismo dos puntos ciegos conocidos, la
distribución de CloudFront, que no participa en las trazas, y el comportamiento
del front-end en el navegador, que no se recoge.

El sprint deja, en definitiva, la monitorización establecida antes de que el
sistema incorpore la persistencia de datos y el procesamiento de recetas. Cada
\textit{endpoint} que se añada en adelante quedará instrumentado sin trabajo
adicional, y las decisiones técnicas de los sprints siguientes podrán apoyarse
en medidas tomadas sobre el sistema en funcionamiento y no únicamente en el
análisis previo.
